\section{The algorithm and its running time}
\label{sec:lvv_query_model}

Theorem~1 of \cite{Landau2015Polynomial} asserts that the ground state is found
in time $n^{c(d,\varepsilon)}\,\mathrm{poly}(n/\eta)$. The paper counts each
linear-algebra step as one polynomial-time computation on data of polynomial
dimension and does not analyse it further. This section fixes the cost model in
which that count is made and states the running time in it; the declared
costs, the honest answering of every operation and exact arithmetic are
recorded in~\cite{gap:lvv13_cost_model}.

A computation is a finite tree of operations: it either returns a value or
performs one operation and continues with a computation depending on the
answer. An oracle assigns an answer to every operation; running a computation
with an oracle produces its value, and the cost of the run is the sum of the
declared costs of the operations performed along the way.

\begin{definition}[The linear-algebra operations]
    \label{def:lvv_linear_algebra_query}
    \lean{TNLean.MPS.AreaLaw.LinearAlgebraQuery,
        TNLean.MPS.AreaLaw.LinearAlgebraQuery.honestOracle,
        TNLean.MPS.AreaLaw.LinearAlgebraQuery.queryCost,
        TNLean.MPS.AreaLaw.LinearAlgebraQuery.one_le_queryCost,
        TNLean.MPS.AreaLaw.extensionCost,
        TNLean.MPS.AreaLaw.netCost,
        TNLean.MPS.AreaLaw.convexProgramCost,
        TNLean.MPS.AreaLaw.schmidtTrimCost,
        TNLean.MPS.AreaLaw.filterFactorsCost,
        TNLean.MPS.AreaLaw.applyFactorCost,
        TNLean.MPS.AreaLaw.energyMinimizerCost}
    \leanok
    \uses{def:lvv_sweep_step, def:lvv_finite_boundary_selection,
        def:chain_boundary_trace_norm_net, def:huang_gaussian_filter_factors,
        def:lvv_energy_minimizer}
    The operations of the algorithm, each with its honest answer and its
    declared cost, are the following. Every operation carries the chain length
    $L$ and the local dimension $q$; the cost is one unit plus a polynomial in
    the dimensions named.
    \begin{enumerate}
        \item The extension $v\otimes|k\rangle$ of a vector $v$ on the first
              $i-1$ sites by a basis state of site $i$, for a vector of bond
              dimension at most $b$: cost $1+qb^2$.
        \item The trace-norm net of radius $\eta$ on the boundary contractions at
              a cut with a bond of dimension $m$; the honest answer is the net
              of Definition~\ref{def:chain_boundary_trace_norm_net}: cost one
              plus $(qm)^2$ times the cardinality bound of the net.
        \item Feasibility of the size-trimming convex program over the span of a
              family $S$ of $s$ vectors of bond dimension at most $b$, at a
              boundary point $X$ with tolerance $\eta$, and the leading
              eigenvector of a solution of least left energy; the honest answers
              are those of Definition~\ref{def:lvv_finite_boundary_selection}:
              cost $1+Lqb^2(sm)^3$ each.
        \item The Schmidt trimming of a vector of bond dimension at most $b$ to
              $m$ terms at every internal bond: cost $1+Lqb^3$.
        \item The left factors at a cut of the Gaussian filters of
              Definition~\ref{def:huang_gaussian_filter_factors}, $E$ filters
              of rank at most $R$: cost $1+ERLq^2$.
        \item The application of a left factor of rank at most $r$ to a vector
              of bond dimension at most $b$: cost $1+Lq^2r^2b^2$.
        \item A unit vector of least energy in the span of a family of $s$
              vectors of bond dimension at most $b$ on the whole chain, the
              final convex program; the honest answer is the vector of
              Definition~\ref{def:lvv_energy_minimizer}: cost $1+Lqb^2s^3$.
    \end{enumerate}
\end{definition}

\begin{definition}[The solution of the final convex program]
    \label{def:lvv_energy_minimizer}
    \lean{TNLean.MPS.AreaLaw.exists_energy_minimizer,
        TNLean.MPS.AreaLaw.energyMinimizer,
        TNLean.MPS.AreaLaw.energyMinimizer_spec}
    \leanok
    Let $H$ be a matrix on a finite-dimensional space and $S$ a finite family
    of vectors whose span $U$ is nonzero. The unit sphere of $U$ is compact and
    the energy $v\mapsto\operatorname{Re}\langle v,Hv\rangle$ is continuous,
    so a unit vector of $U$ of least energy exists; one such vector is chosen,
    and $0$ is taken when the span is zero. It is the leading eigenvector of a
    pure optimal solution of the final convex program of
    \cite{Landau2015Polynomial}.
\end{definition}

\begin{definition}[The steps of the algorithm as computations]
    \label{def:lvv_step_programs}
    \lean{TNLean.MPS.AreaLaw.LinearAlgebraQuery.unionOver,
        TNLean.MPS.AreaLaw.LinearAlgebraQuery.eval_unionOver,
        TNLean.MPS.AreaLaw.LinearAlgebraQuery.cost_unionOver,
        TNLean.MPS.AreaLaw.leftColumnsOf,
        TNLean.MPS.AreaLaw.extensionProgram,
        TNLean.MPS.AreaLaw.sizeTrimmingProgram,
        TNLean.MPS.AreaLaw.bondTrimmingProgram,
        TNLean.MPS.AreaLaw.errorReductionProgram,
        TNLean.MPS.AreaLaw.sweepStepProgram,
        TNLean.MPS.AreaLaw.sweepProgram}
    \leanok
    \uses{def:lvv_linear_algebra_query, def:lvv_sweep}
    A computation on a finite family performs a computation at each member and
    unites the finite sets returned. The four steps of the algorithm are the
    computations that extend every member of the family by every basis state;
    for every bond dimension up to the uniform cut bound enumerate the net,
    test feasibility at every point and keep the left columns of the leading
    eigenvector at every feasible point; trim every member; and apply every
    left factor of the Gaussian filters to every member. One step of the sweep
    performs them in order with the bond budget of
    Definition~\ref{def:lvv_successive_budgets} declared for the vectors
    handled, and the sweep performs the step at every cut starting from the
    scalar singleton.
\end{definition}

\begin{lemma}[The computations compute the sweep]
    \label{thm:lvv_eval_sweep_program}
    \lean{TNLean.MPS.AreaLaw.eval_extensionProgram,
        TNLean.MPS.AreaLaw.eval_sizeTrimmingProgram,
        TNLean.MPS.AreaLaw.eval_bondTrimmingProgram,
        TNLean.MPS.AreaLaw.eval_errorReductionProgram,
        TNLean.MPS.AreaLaw.eval_sweepStepProgram,
        TNLean.MPS.AreaLaw.eval_sweepProgram}
    \leanok
    \uses{def:lvv_step_programs}
    Run with the honest oracle, the four computations return the extended
    family, the size-trimming selection, the trimmed family and the
    error-reduced family of the mathematical development, the step returns the
    step of the sweep, and the sweep computation returns the sweep at every
    cut. In particular the outputs satisfy the viability, cardinality and bond
    bounds of Lemma~\ref{thm:lvv_sweep_spec}.
\end{lemma}
\begin{proof}
    \leanok
    Running a union over a family gives the union of the runs, and each run of
    one operation is its honest answer. The identifications are then the
    definitions of the four families, read member by member; the size-trimming
    selection is the union over bond dimensions of the left columns at the
    feasible net points. The step and the sweep follow by composition and by
    induction on the cut.
\end{proof}

\begin{definition}[The ground-state computation]
    \label{def:lvv_ground_state_program}
    \lean{TNLean.MPS.AreaLaw.groundStateProgram,
        TNLean.MPS.AreaLaw.eval_groundStateProgram}
    \leanok
    \uses{def:lvv_step_programs, def:lvv_energy_minimizer,
        def:lvv_padded_interaction}
    For an accuracy parameter $R\ge1$, pad the chain by $2R^2$ sites with the
    penalty of Definition~\ref{def:lvv_padded_interaction}, run the sweep
    computation on the padded chain at gap parameter $\min\{\Delta,1\}$ up to
    the cut after the original $L$ sites, read the final family on the original
    chain, and return the unit vector of least energy in its span. Run with the
    honest oracle, the output is the chosen energy minimizer over the final
    family of the padded sweep.
\end{definition}

\begin{theorem}[The ground-state computation finds the ground state]
    \label{thm:lvv_ground_state_program_spec}
    \lean{TNLean.MPS.AreaLaw.one_sub_le_norm_inner_of_norm_sub_sq_le,
        TNLean.MPS.AreaLaw.groundStateProgram_spec}
    \leanok
    \uses{def:lvv_ground_state_program, thm:lvv_sweep_spec,
        thm:lvv_padding_prefix_mps, thm:chain_range_one_overlap_energy,
        thm:lvv_left_span_bond}
    Let $L\ge1$, $q\ge2$, $\Delta>0$ and $R\ge1$, and let $H=\sum_j h_j$ be an
    adjacent-pair interaction with $0\le h_j\le I$. For every unit ground
    vector $\Gamma$ with projector gap $\Delta$, the output $\psi$ of the
    ground-state computation with the honest oracle is a unit vector with
    \begin{align}
        |\langle\Gamma,\psi\rangle|\ge1-\frac{4L}{R^2\Delta},
    \end{align}
    and an open-boundary matrix product state whose bond dimension is at most
    the cardinality budget times the bond budget of the sweep on the padded
    chain. Choosing $R^2\ge4L/(\Delta\eta)$ gives overlap at least $1-\eta$.
\end{theorem}
\begin{proof}
    \leanok
    By Lemma~\ref{thm:lvv_padding_prefix_mps} the span of the final family of
    the padded sweep, read on the original chain, contains a unit vector $w$
    with $\|w-\Gamma\|^2\le2\delta$, where $\delta=c_{\min\{\Delta,1\}}/N$ and
    $N=L+2R^2$; a unit vector within squared distance $2\delta$ of a unit
    vector has overlap at least $1-\delta$ with it. By
    Lemma~\ref{thm:chain_range_one_overlap_energy} the energy of $w$ exceeds the
    ground energy by at most $8L\delta$, hence so does the least energy of the
    span, and the gap lemma gives overlap at least $1-8L\delta/\Delta$ for the
    minimizer; since $\delta\le1/(2R^2)$ this is at least $1-4L/(R^2\Delta)$.
    The bond bound is that of a vector in the span of the final family.
\end{proof}

\begin{definition}[Polynomially bounded functions]
    \label{def:nat_polynomial_bounded}
    \lean{TNLean.IsNatPolynomialBounded,
        TNLean.IsNatPolynomialBounded.const,
        TNLean.IsNatPolynomialBounded.id,
        TNLean.IsNatPolynomialBounded.of_le,
        TNLean.IsNatPolynomialBounded.add,
        TNLean.IsNatPolynomialBounded.mul,
        TNLean.IsNatPolynomialBounded.pow,
        TNLean.IsNatPolynomialBounded.max_one,
        TNLean.IsNatPolynomialBounded.exists_one_le}
    \leanok
    A function $f$ from the natural numbers to themselves is polynomially
    bounded when $f(n)\le P\,(n+1)^d$ for all $n\ge1$, for natural numbers $P$
    and $d$ chosen before $n$. Constants and the identity are polynomially
    bounded, and the class is closed under pointwise domination, sums,
    products, powers and taking the maximum with one; the coefficient $P$ can be
    taken positive.
\end{definition}

\begin{definition}[The cost budgets]
    \label{def:lvv_cost_budgets}
    \lean{TNLean.MPS.AreaLaw.sweepStepCostBound,
        TNLean.MPS.AreaLaw.sweepStepCostBound_mono,
        TNLean.MPS.AreaLaw.groundStateCostBound,
        TNLean.MPS.AreaLaw.netCost_mono,
        TNLean.MPS.AreaLaw.convexProgramCost_mono,
        TNLean.MPS.AreaLaw.energyMinimizerCost_mono}
    \leanok
    \uses{def:lvv_linear_algebra_query, def:lvv_successive_budgets}
    The cost budget of one step of the sweep on a family of at most $s$
    members is the sum of the declared costs of $qs$ extensions, of one net
    enumeration and at most two convex programs per net point for every bond
    dimension up to the uniform cut bound, of one trimming per member of the
    selection, of the filter factors, and of one application per factor and
    trimmed member, with the cardinality and bond budgets of
    Definition~\ref{def:lvv_successive_budgets} in place of the actual
    dimensions. The budget of the ground-state computation on the padded chain
    of $N$ sites is $N$ steps at the cardinality budget plus the final convex
    program.
\end{definition}

\begin{lemma}[The cost of the computations]
    \label{thm:lvv_program_cost}
    \lean{TNLean.MPS.AreaLaw.cost_extensionProgram_le,
        TNLean.MPS.AreaLaw.cost_feasible_ite_le,
        TNLean.MPS.AreaLaw.cost_sizeTrimmingProgram_le,
        TNLean.MPS.AreaLaw.cost_bondTrimmingProgram_le,
        TNLean.MPS.AreaLaw.cost_errorReductionProgram_le,
        TNLean.MPS.AreaLaw.cost_sweepStepProgram_le,
        TNLean.MPS.AreaLaw.cost_sweepProgram_le,
        TNLean.MPS.AreaLaw.cost_groundStateProgram_le}
    \leanok
    \uses{def:lvv_cost_budgets, thm:lvv_eval_sweep_program,
        thm:lvv_size_trimmed_family_polynomial, def:chain_boundary_trace_norm_net}
    Run with the honest oracle, each step costs at most its budget, one step of
    the sweep on a family $S$ costs at most the step budget at $|S|$, the sweep
    up to cut $i\le L$ costs at most $i$ step budgets at the cardinality budget,
    and the ground-state computation costs at most its budget on the padded
    chain.
\end{lemma}
\begin{proof}
    \leanok
    The cost of a union over a family is the sum of the costs, bounded by the
    cardinality times the largest cost. The extended family has at most $q|S|$
    members, the net has the cardinality bound of
    Definition~\ref{def:chain_boundary_trace_norm_net}, the selection has the polynomial
    cardinality of Lemma~\ref{thm:lvv_size_trimmed_family_polynomial}, trimming
    does not increase cardinality, and the filter factors number at most the
    energy count times the rank bound. The families of the sweep have at most
    the cardinality budget members by Lemma~\ref{thm:lvv_sweep_spec}, which
    bounds every step uniformly, and the number of steps is at most the padded
    chain length.
\end{proof}

\begin{theorem}[Polynomial running time]
    \label{thm:lvv_polynomial_running_time}
    \lean{TNLean.MPS.AreaLaw.isNatPolynomialBounded_sweepStepCostBound,
        TNLean.MPS.AreaLaw.isNatPolynomialBounded_groundStateCostBound,
        TNLean.MPS.AreaLaw.exists_groundStateProgram_cost_polynomial}
    \leanok
    \uses{thm:lvv_program_cost, thm:lvv_successive_polynomial_budgets,
        lem:lvv_bond_trimming_radius_polynomial, def:nat_polynomial_bounded}
    Fix $q\ge2$ and $\Delta>0$. There are natural numbers $P\ge1$ and $d$,
    chosen before the chain, such that for every chain of $L\ge1$ sites, every
    adjacent-pair interaction with $0\le h_j\le I$ and every accuracy parameter
    $R$, the ground-state computation run with the honest oracle performs at
    most as many operations as its cost, and its cost is at most
    $P\,(L+2R^2+1)^d$. With $R^2$ of order $L/(\Delta\eta)$ this is the bound
    $n^{c}\,\mathrm{poly}(n/\eta)$ of Theorem~1 of
    \cite{Landau2015Polynomial}, with $c$ depending on $q$ and $\Delta$ through
    $P$ and $d$; the printed form $c=2^{O(\log^3 d/\varepsilon)}$ of the exponent
    is not reproduced, as explained in~\cite{gap:lvv13_cost_model}.
\end{theorem}
\begin{proof}
    \leanok
    Every quantity entering the step budget is bounded by a polynomial in the
    chain length with coefficients chosen before it: the cardinality and bond
    budgets and the filter counts by
    Lemma~\ref{thm:lvv_successive_polynomial_budgets}, the net cardinality and
    the net cost by the corrected net exponent, the selection cardinality by its
    definition, and the trimming radius by
    Lemma~\ref{lem:lvv_bond_trimming_radius_polynomial}. Polynomially bounded
    functions are closed under sums, products and powers, so the step budget,
    and then the budget of the ground-state computation on the padded chain, are
    polynomially bounded. Every operation costs at least one unit, so the
    number of operations is at most the cost.
\end{proof}
